\documentclass[10pt,a4paper]{amsart}
\usepackage[margin=19mm]{geometry}
\usepackage{amsmath,amssymb,mathtools}
\usepackage[scr=rsfs,cal=euler]{mathalfa}
\usepackage{xfrac}
\usepackage[unicode,hidelinks,bookmarks=false]{hyperref}
\newcommand{\ie}{i.e.\ }
\newcommand{\cf}{cf.\ }
\newcommand{\resp}{respectively\ }
\newcommand{\Subsection}{\subsection}
\providecommand{\nobreakdash}{\nobreak\hbox{-}}
\setlength{\emergencystretch}{2em}
\multlinegap0pt
\usepackage{bm}
\usepackage{bbm}
\usepackage{dsfont}
\usepackage{xspace}
\usepackage{cancel}
\usepackage{mathtools}
\usepackage{amsthm}
\makeatletter
\@ifundefined{definition}{\newtheorem{definition}{Definition}}{}
\@ifundefined{lemma}{\newtheorem{lemma}[definition]{Lemma}}{}
\makeatother
\renewcommand{\geq}{\geqslant}
\renewcommand{\ge}{\geqslant}
\renewcommand{\leq}{\leqslant}
\renewcommand{\le}{\leqslant}
\title[A free boundary analysis of tumor invasion driven by angioge]{A free boundary analysis of tumor invasion driven by angiogenesis}
\author{Serena Dipierro, Edoardo Proietti Lippi, Caterina Sportelli, Enrico Valdinoci}
\date{Source: arXiv:2607.05610v2 (CC BY 4.0)}
\begin{document}
\maketitle

\noindent\emph{Source and licence.} A free boundary analysis of tumor invasion driven by angiogenesis. Version arXiv:2607.05610v2 (CC BY 4.0), distributed under the Creative Commons Attribution 4.0 International licence, as recorded in the arXiv licence metadata for this exact version and shown on the abstract page https://arxiv.org/abs/2607.05610v2. Full rights notice: licenses/arxiv-2607.05610-CC-BY-4.0.txt.\par\smallskip

\noindent\textit{Editorial scope.} Counted unit: the Lemma whose source label is \texttt{lemmastimespunto} in the article (source0.tex lines 1080--1092, source label \texttt{lemmastimespunto}), delivered together with the article's own complete original proof of it (source0.tex lines 1094--1143, one \texttt{proof} environment of 50 lines). No mathematics is written, repaired or re-derived: statement and proof are the article's own text line for line. Required context retained uncounted: u\_0positiva (lines 301--303); mainfinal1 (lines 315--324); freeboundary (lines 425--427); relationutilde2 (lines 445--447); lemma0<u<utilde (lines 1034--1046). Every definition, display and label of those lines is retained unchanged. Omitted from this package and not redistributed: 1--300, 304--314, 325--424, 428--444, 448--1033, 1047--1079, 1093--1093, 1144--3407. Nothing inside a retained excerpt is omitted or altered: every retained body is delivered exactly as the article's lines are written, so no omission receipt of any kind accompanies this package. Citations: the retained excerpts carry no citation command at all, so no bibliography entry of the article is transcribed and no reference is redistributed. Reference adaptation: none was needed and none was made; every reference inside the delivered unit resolves inside it, so no unresolved reference or citation is produced. Printed slips kept literal: no printed word, symbol, hypothesis or number of the counted statement or of its proof was altered. Layout and class: the article's own class is \texttt{amsart} with packages including amsmath, amsthm, mathtools, amsfonts, amsrefs, graphicx, framed, amssymb, esvect, xcolor, eucal, mathrsfs, hyperref, babel; the wrapper is the lane's pinned \texttt{amsart} editorial wrapper at 10pt on A4, which restates the theorem-like environments the retained text needs and adds the article's own macro definitions that the retained text needs (\texttt{\textbackslash ge}, \texttt{\textbackslash geq}, \texttt{\textbackslash le}, \texttt{\textbackslash leq}), with extra wrapper packages (bm, bbm, dsfont, xspace, cancel, mathtools). The delivered numbering is the unit's own. Assets: no retained span contains a figure environment, a graphics-inclusion command, a PDF-page inclusion, a table or a TikZ picture; no image file is copied into this package, the package declares no asset dependency and no image file is redistributed with it. Licence: version arXiv:2607.05610v2 of this article is distributed under the Creative Commons Attribution 4.0 International licence, as recorded in the arXiv licence metadata for this exact version and shown on the abstract page https://arxiv.org/abs/2607.05610v2. A full rights notice is delivered at licenses/arxiv-2607.05610-CC-BY-4.0.txt. Characters: every retained excerpt is pure ASCII, so the delivered engine, XeLaTeX with its default Latin Modern text font, reports no missing character.
\begin{equation}\label{u_0positiva}
0<\frac{\Gamma u_B}{\lambda}\le u_0(r)\le\overline{u}, \qquad \mbox{if }\,\,\varepsilon\leq r\leq s(0).
\end{equation}

\begin{equation}\label{mainfinal1}
\begin{cases}
\kappa u_t (x,t)= \Delta u(x, t) - \displaystyle\frac{\chi_0}{r^2}\, x\cdot\nabla u(x, t) -\lambda u(x, t)+ \Gamma u_B 
%-\lambda u(x, t) 
&\quad\mbox{ in }B_{s(t)}\setminus\overline{B_\varepsilon},\\
u(x, t) = \underline{u}, &\quad\mbox{ on } \partial B_\varepsilon \\
u(x, t) = \overline{u} &\quad\mbox{ on } \partial B_{s(t)}, \\
u(x, 0) = u_0(|x|) &\quad\mbox{ in } B_{s(0)}\setminus B_\varepsilon.
\end{cases}
\end{equation}

\begin{equation}\label{freeboundary}
s^2(t)\dot s(t)=\mu\int_\varepsilon^{\max\{s(t),\varepsilon\}}\big(u(r,t)-\widetilde u\big)r^2\,dr.
\end{equation}

\begin{equation}\label{relationutilde2}
0<\frac{\Gamma u_B}{\lambda}<\widetilde u < \underline{u} < \overline{u}.
\end{equation}

\begin{lemma}\label{lemma0<u<utilde}
Assume~\eqref{u_0positiva} and~\eqref{relationutilde2}.
Let $(u(r,t), s(t))$ be a solution of~\eqref{mainfinal1} and~\eqref{freeboundary}.

Then, for all $\varepsilon\le r\le s(t)$ and~$t\ge0$,
\begin{equation}\label{prima0}
u(r,t)\ge \frac{\Gamma u_B}{\lambda}>0
\end{equation}
and
\begin{equation}\label{prima4}
u(r,t)\leq\overline{u}.
\end{equation}
\end{lemma}

\begin{lemma}\label{lemmastimespunto}
Assume~\eqref{u_0positiva} and~\eqref{relationutilde2}.
Let $(u(r,t), s(t))$ be a solution of~\eqref{mainfinal1} and~\eqref{freeboundary}.

Then, for all $t>0$,
\begin{equation}\label{stimespunto}
-\mu\left(\widetilde u-\frac{\Gamma u_B}{\lambda}\right)(s(t)-\varepsilon)\leq \dot s(t)\leq \mu(\overline{u}-\widetilde u)(s(t)-\varepsilon)
\end{equation}
and
\begin{equation}\label{stimespunto2}
(s(0)-\varepsilon)e^{-\mu\left(\widetilde u-\frac{\Gamma u_B}{\lambda}\right)t} +\varepsilon\leq s(t)\leq (s(0)-\varepsilon)e^{\mu(\overline{u}-\widetilde u) t}+\varepsilon .
\end{equation}
\end{lemma}

\begin{proof} We set $$
T:=\sup\{ t\in[0,+\infty) {\mbox{ s.t. }} s(t)>\varepsilon \}$$
and we remark that~$T\in(0,+\infty]$, since~$s(0)>\varepsilon$.

Our strategy is to establish~\eqref{stimespunto}
and~\eqref{stimespunto2} for all~$t\in[0,T)$. If we accomplish this,
then the first inequality in~\eqref{stimespunto2} yields that~$T=+\infty$,
hence giving~\eqref{stimespunto}
and~\eqref{stimespunto2} for all~$t\in[0,+\infty)$.

This approach is helpful, because for all~$t\in[0,T)$
the free boundary condition~\eqref{freeboundary} boils down to
\begin{equation}\label{QUEST}s^2(t)\dot s(t)=\mu\int_\varepsilon^{s(t)}\big(u(r,t)-\widetilde u\big)r^2\,dr.\end{equation}

Now,
thanks to Lemma~\ref{lemma0<u<utilde} we know that~$u\ge\frac{\Gamma u_B}{\lambda}$. Thus,
from~\eqref{QUEST} we deduce that, for all~$t\in[0,T)$,
\begin{equation*}\begin{split}&
\dot{s}(t)
=\frac{\mu}{s^2(t)}\int_\varepsilon^{s(t)}\big(u(r,t)-\widetilde u\big)r^2\,dr
\geq-\frac{\mu\left(\widetilde u-\frac{\Gamma u_B}{\lambda}\right)}{s^2(t)}\int_\varepsilon^{s(t)}r^2\,dr=
-\frac{\mu\left(\widetilde u-\frac{\Gamma u_B}{\lambda}\right)(s^3(t)-\varepsilon^3)}{3s^2(t)}\\&\qquad
=
-\mu\left(\widetilde u-\frac{\Gamma u_B}{\lambda}\right)(s(t)-\varepsilon)\frac{s^2(t)+\varepsilon s(t)+\varepsilon^2}{3s^2(t)}.\end{split}
\end{equation*}
Since~$\varepsilon<s(t)$, we have that
$$ \frac{s^2(t)+\varepsilon s(t)+\varepsilon^2}{3s^2(t)}=
\frac13\left(1+\frac{\varepsilon}{s(t)}+\frac{\varepsilon^2}{s^2(t)}\right)\le 1
$$ and therefore, for all~$t\in[0,T)$,
\begin{equation}\label{zxdrtfvghyuhnj1} \dot{s}(t)
\ge -\mu\left(\widetilde u-\frac{\Gamma u_B}{\lambda}\right)(s(t)-\varepsilon). \end{equation}

We observe that Lemma~\ref{lemma0<u<utilde} also tells us that~$u\le \overline{u}$. As a result,
using this information into~\eqref{QUEST} we get that, for all~$t\in[0,T)$,
\begin{equation*}
\begin{split}
&\dot{s}(t)\leq \frac{\mu(\overline{u}-\widetilde u)}{s^2(t)}\int_\varepsilon^{s(t)}r^2\,dr
=\frac{\mu(\overline{u}-\widetilde u)(s^3(t)-\varepsilon^3)}{3s^2(t)}
\\&\qquad=\frac{\mu}{3}(\overline{u}-\widetilde u)(s(t)-\varepsilon)\frac{s^2(t)+\varepsilon s(t)+\varepsilon^2}{s^2(t)}
\le \mu(\overline{u}-\widetilde u)(s(t)-\varepsilon).
\end{split}
\end{equation*}
Combining this and~\eqref{zxdrtfvghyuhnj1}, we obtain the desired result in~\eqref{stimespunto}, for all~$t\in[0,T)$.

Furthermore, by integrating~\eqref{stimespunto},  we deduce that
\[
-\mu\left(\widetilde u-\frac{\Gamma u_B}{\lambda}\right) t\leq \log(s(t)-\varepsilon)-\log(s(0)-\varepsilon)\leq \mu(\overline{u}-\widetilde u) t,
\]
which entails~\eqref{stimespunto2}, for all~$t\in[0,T)$, as desired.
\end{proof}

\end{document}
